\begin{tikzpicture}[
    every node/.style={treeblue},
    line/.style={treeblue, line width=1pt},
]
% Enrique Acosta, 2026
% ============================================================================
% A full n-ary tree drawn sideways: root on the left, leaves fanning out on
% the right. 
% 
% Growth direction is controlled purely by the coordinate formula. 
% Swap \xx and \yy in (1) to get a tree opening upwards, add a "-" to 
% get it to open downwards.
% ============================================================================

% --- Tree parameters -----------------------------------------------
\def\branch{3}       % children per node (branching factor)
\def\depth{4}        % number of splits = tree depth (there are \depth+1 rows)
% --- Layout parameters ------------------------------
\def\levelsep{2.2}   % gap between levels (cm) in the branching direction
\def\leafsep{0.2}    % perpendicular gap between adjacent leaves (cm)
\def\rootrad{3.2}    % dot radius at the root (pt)
\def\leafrad{1.6}    % dot radius at the leaves (pt), interpolated in between
% -------------------------------------------------------------

% --- Derived constants (computed once from the Tree parameters) -------- 
\pgfmathtruncatemacro{\nleaves}{\branch^\depth}    % total leaves
\pgfmathsetmacro{\midleaf}{(\nleaves-1)/2}         % center leaf index

% 1) Define a coordinate for every node.
%    Depth d (0..\depth) has branch^d nodes; each is centered on the leaves
%    it covers.
\foreach \d in {0,...,\depth} {
    \pgfmathtruncatemacro{\nmax}{\branch^\d - 1}         % last node index here
    \pgfmathtruncatemacro{\leavesper}{\branch^(\depth-\d)}% leaves under one node
    \foreach \i in {0,...,\nmax} {
%       Midpoint of the leaf block [ \i*\leavesper , +\leavesper-1 ], in
%       leaf-index units. Centering each node on its block keeps parents
%       symmetric above their children.
        \pgfmathsetmacro{\cen}{\i*\leavesper + (\leavesper-1)/2}
%       Recenter on 0 (subtract the middle leaf index) and scale to cm.
        \pgfmathsetmacro{\yy}{(\cen-\midleaf)*\leafsep}
        \pgfmathsetmacro{\xx}{\d*\levelsep}
        \coordinate (n-\d-\i) at (\xx,\yy);
    }
}

% 2) Draw the edges. Loop over PARENT depths 0..\depth-1; each parent \i has
%    children \branch*\i + \k for \k = 0..\bm on the next depth.
%    Depth=0 is special: has arrowheads and space between them an the child nodes
\pgfmathtruncatemacro{\dm}{\depth-1}  % last PARENT depth
\pgfmathtruncatemacro{\bm}{\branch-1} % last child offset
\foreach \d in {0,...,\dm} {
    \pgfmathtruncatemacro{\nmax}{\branch^\d - 1}
    \pgfmathtruncatemacro{\dn}{\d+1}
    \foreach \i in {0,...,\nmax} {
        \foreach \k in {0,...,\bm} {
            \pgfmathtruncatemacro{\child}{\branch*\i+\k}
            \ifnum\d=0
                \draw[line, -{Stealth[length=3mm]}, shorten >=4pt]
                    (n-0-\i) -- (n-\dn-\child);
            \else
                \draw[line] (n-\d-\i) -- (n-\dn-\child);
            \fi
        }
    }
}

% 3) Draw the dots on top, radius interpolated from \rootrad to \leafrad
%    linearly with depth (replaces the old hand-picked size thresholds).
\foreach \d in {0,...,\depth} {
    \pgfmathtruncatemacro{\nmax}{\branch^\d - 1}
    \pgfmathsetmacro{\rad}{\rootrad + (\leafrad-\rootrad)*\d/\depth}
    \foreach \i in {0,...,\nmax} {
        \fill[treeblue] (n-\d-\i) circle (\rad pt);
    }
}

% 4) Root label and annotation.
\node[ anchor=south
     % , font=\large
     , text width=1.5cm
     , align=center
     %, draw
     ] at ([xshift=-4mm, yshift=6mm]n-0-0) {No choice yet!};

\node[ anchor= south
     % , font=\itshape
     , text width=4cm
     , align=center
     %, draw
     ] at ([xshift=-10mm, yshift=4mm]n-1-2)
     {This tree splits in three four times.\\ It has 4 levels.\\
      It shows all 81 músicos.\\($3\times3\times3\times3$)};

\end{tikzpicture}